\scanpage{29}\mspage{28}
Soient $F$ et $\unsure{G|C}$ deux figures. Conditions équivalentes (car $F$ et $G$ sont des ens. \uline{admissibles} de parties de $\mathcal{L}$)

a) $|F|\cap|G|$ est une partie saturée de $|F|$ et de $|G|$, et les deux rel, d'équiv. induites sont égales

% slanted
\margin{\uline{NB} $F$ et $G$ ens\textsuperscript{t} comp. sss \deleted{toute} \deleted{strate}\inserted{multistrate} de $F$ est ens. comp. à \deleted{toute} multistrate de $G$}

b) $\forall$ strate ouverte $U$\inserted{$=X^\circ$} de $F$ et strate ouverte \deleted{$\unsure{U|V}$}$V$\inserted{$={\unsure{Y|Z}}^\circ$} de $G$, on a $U\cap V=\varnothing$ ou $U=V$

b$'$) $\ldotp\ldotp\ldotp\ldotp\ldotp$ \deleted{On} a $\deleted{\unsure{U}}X^\circ\cap\deleted{\unsure{V}}Y^\circ=\varnothing$ ou $X=Y$

c) $\forall\,x\in F,Y\in G$, $X\cap Y$ est saturé dans $|F|$ et de $|G|$. \inserted{(réentraine la saturation pour les str. des multistr. \deleted{qui} $\widehat{x},\widehat{y}$)}

d) $F\cup G$ est admissible,) et $F$ et $G$ sont des sous-ens\unsure{.} \deleted{fer} \uline{fermés} pour la relation $\subset$

On dit alors que \deleted{\unsure{$G$,}} $F$ et $G$ sont ``ensemblistement compatibles''. Il en résulte que

% slanted
\margin{\uline{NB} Soit $H$ une autre famille admissible d'ens. Pour que $H$ soit ens. comp. à $F$ et $G$ (supposés ens. comp. entre eux), il f et s qu'il le soit à $F\cup G$.}
\[
\begin{aligned}
&\relax\inserted{F\cap G\ \text{fermé dans } F \text{ et dans } G}\\
&F,G\ \text{fermés dans } F\cup G
\end{aligned}
\]

et $F\cup G$, en tant qu'ens \inserted{(?)} ordonné est l'\deleted{\unsure{ens.}} \deleted{\unsure{som}} somme \unsure{amalgamée} $>\quad F\amalg_{F\cap G}G$,

% slanted
\margin{\uline{NB} Si $F$ et $G$ sont compatibles alors deux sous figures $F'\leq F$ et $G'\leq G$ \unsure{le sont ssi|le sont ss.|le sont aussi} Donc si on a une configuration, l'ens des sous-figures de toutes les figures de la conf. est encore une configuration. On peut aussi \unsure{agrandir|adj.} en y incluant toutes les réunions finies \unsure{d'ens.|de} \deleted{telles} figures \unsure{de la conf.(ig)|de la config} $+$ arbitraires.}

On dit que $F$ et $G$ \uline{sont compatibles}, \deleted{si $F\cup G$ est} s'ils le sont ensemblis\deleted{\unsure{st|t}}tement, et si $F\cup G$ est une figure. \inserted{(Alors $F,G$ sont des sous-figures de $F\cup\unsure{C|G}$)}

\uline{Ex} Si $F$ et $G$ sont des sous-figures d'une figure $H$, alors ils sont compatibles (car \inserted{ils le sont ensembl\textsuperscript{t}., et} $F\cup G$ est une sous-figure de $H$, donc une figure.)

Configuration : ensemble de figures deux à deux compatibles.
